%-------------------------
% Resume in LateX
% Author : Gabriel Gomez
% License : MIT
%------------------------

\documentclass[letterpaper,11pt]{article}

\usepackage{cmap}
\usepackage[utf8]{inputenc}
\usepackage[T1]{fontenc}
\usepackage{lmodern}

\usepackage{latexsym}
\usepackage[empty]{fullpage}
\usepackage{titlesec}
\usepackage{marvosym}
\usepackage[usenames,dvipsnames]{color}
\usepackage{verbatim}
\usepackage{enumitem}
\usepackage[hidelinks]{hyperref}
\usepackage{fancyhdr}
\usepackage[spanish]{babel}
\usepackage{tabularx}
\input{glyphtounicode}

\pagestyle{fancy}
\fancyhf{} % Clear all header and footer fields
\fancyfoot{}
\renewcommand{\headrulewidth}{0pt}
\renewcommand{\footrulewidth}{0pt}

% Adjust margins
\addtolength{\oddsidemargin}{-0.5in}
\addtolength{\evensidemargin}{-0.5in}
\addtolength{\textwidth}{1in}
\addtolength{\topmargin}{-.5in}
\addtolength{\textheight}{1.0in}

\urlstyle{same}

\raggedbottom
\raggedright
\setlength{\tabcolsep}{0in}

% Sections formatting
\titleformat{\section}{
  \vspace{-4pt}\scshape\raggedright\large
}{}{0em}{}[\color{black}\titlerule \vspace{-5pt}]

% Ensure that generate PDF is machine readable/ATS parsable
\pdfgentounicode=1

%-------------------------
% Custom commands
\newcommand{\resumeItem}[2]{
  \item\small{
    \textbf{#1}{: #2 \vspace{-2pt}}
  }
}

% Just in case someone needs a heading that does not need to be in a list
\newcommand{\resumeHeading}[4]{
    \begin{tabular*}{0.99\textwidth}[t]{l@{\extracolsep{\fill}}r}
      \textbf{#1} & #2 \\
      \textit{\small#3} & \textit{\small #4} \\
    \end{tabular*}\vspace{-5pt}
}

\newcommand{\resumeSubheading}[4]{
  \vspace{-1pt}\item
    \begin{tabular*}{0.97\textwidth}[t]{l@{\extracolsep{\fill}}r}
      \textbf{#1} & #2 \\
      \textit{\small#3} & \textit{\small #4} \\
    \end{tabular*}\vspace{-5pt}
}

\newcommand{\resumeSubSubheading}[2]{
    \begin{tabular*}{0.97\textwidth}{l@{\extracolsep{\fill}}r}
      \textit{\small#1} & \textit{\small #2} \\
    \end{tabular*}\vspace{-5pt}
}

\newcommand{\resumeSubItem}[2]{\resumeItem{#1}{#2}\vspace{-4pt}}

\renewcommand{\labelitemii}{$\circ$}

\newcommand{\resumeSubHeadingListStart}{\begin{itemize}[leftmargin=*]}
\newcommand{\resumeSubHeadingListEnd}{\end{itemize}}
\newcommand{\resumeItemListStart}{\begin{itemize}}
\newcommand{\resumeItemListEnd}{\end{itemize}\vspace{-5pt}}
\setlength{\footskip}{5pt}
%-------------------------------------------
%%%%%%  CV STARTS HERE  %%%%%%%%%%%%%%%%%%%%%%%%%%%%


\begin{document}

%----------HEADING-----------------
\begin{tabular*}{\textwidth}{l@{\extracolsep{\fill}}r}
  \textbf{\href{https://gabogomez.dev/}{\Large Gabriel Gomez}} & Correo electrónico: \href{mailto:ggomez.dev@outlook.com}{ggomez.dev@outlook.com}\\
  \href{https://gabogomez.dev/}{gabogomez.dev} & Movil: \href{tel:+573168325201}{+57-316-832-5201} \\
\end{tabular*}





%-----------EXPERIENCE-----------------
\section{Experiencia}
\resumeSubHeadingListStart

\resumeSubheading
{Autónomo}{Barranquilla, CO}
{Desarrollador de Software Freelance}{Jun 2025 -- Presente}
\resumeItemListStart
\resumeItem{Migración de Web Service}
{Migré un web service de Java EE a Spring Boot e implementé Kubernetes para la orquestación de contenedores. Construí un entorno de desarrollo automatizado usando un Makefile para simplificar flujos de trabajo, incluyendo levantar un stack de Docker Compose con una base de datos de pruebas, construir imágenes de contenedor, correr el proyecto con variables de entorno e inicializar archivos de configuración.}
\resumeItemListEnd

\resumeSubheading
{Extreme Technologies S.A.}{Barranquilla, CO}
{Desarrollador de Software}{Abr 2021 -- Jun 2024}
\resumeItemListStart
\resumeItem{Acuacar Daños}
{Mantuve y desarrollé features para la app Android y el web service Jakarta de un sistema de gestión de cuadrillas en la nube para la empresa de acueducto Acuacar. Rediseñé la interfaz gráfica completa para una apariencia moderna, migré la app para soportar Android 11 e implementé una modificación al flujo de trabajo que abarcó tanto la app móvil como el web service.}
\resumeItem{Promigas}
{Adapté un sistema white-label de gestión de cuadrillas para un nuevo cliente de gas, desarrollando la app Android y el web service Jakarta. Implementé rastreo GPS en tiempo real para mapear las rutas de los contratistas durante inspecciones de gasoductos, integré PostGIS para el procesamiento de coordenadas geográficas y optimicé el web service para manejar la carga incrementada por sesiones concurrentes en tiempo real.}
\resumeItem{Podas y Lavados}
{Adapté un sistema white-label para mantenimiento de redes eléctricas y poda de árboles sirviendo a Aire y Afinia. Construí una app Android multirol con soporte para dos flujos de trabajo paralelos e implementé un feature GIS que descarga y renderiza capas en el mapa con las posiciones y estados de árboles y postes eléctricos desde el web service.}
\resumeItem{SETP Santa Marta}
{Implementé una capa middleware en un validador de pasajes basado en Android para un sistema de transporte público. Construí la validación de pagos para tarjetas físicas y códigos QR de billeteras virtuales, y diseñé una cola offline-first para almacenar transacciones durante pérdidas de conexión y sincronizarlas al restaurar la conectividad.}

\resumeItemListEnd

% --------Multiple Positions Heading------------
%  \resumeSubSubheading
%   {Software Engineer I}{Oct 2014 -- Sep 2016}
%   \resumeItemListStart
%      \resumeItem{Apache Beam}
%        {Apache Beam is a unified model for defining both batch and streaming data-parallel processing pipelines}
%   \resumeItemListEnd

%-------------------------------------------

\resumeSubHeadingListEnd

%-----------EDUCATION-----------------
\section{Educación}
\resumeSubHeadingListStart
\resumeSubheading
{Universidad Autónoma del Caribe}{Barranquilla, CO}
{Pregrado en Ingeniería de sistemas}{Feb 2018 -- May 2023}
\resumeSubHeadingListEnd


%-----------PROJECTS-----------------
\section{Proyectos}
\resumeSubHeadingListStart
\resumeSubItem{Weather App}
{Aplicación meteorológica de código abierto para Android creada con Kotlin y Jetpack Compose que ofrece previsiones en tiempo real, actualizaciones basadas en la ubicación y una interfaz de usuario moderna y personalizable mediante la API OpenWeather.}
\resumeSubItem{Resume Builder on LateX}
{Un generador de currículos de código abierto que utiliza LaTeX para crear PDF de aspecto profesional y utiliza GitHub Actions para la creación y el despliegue automatizados.}
\resumeSubItem{Portfolio Personal}
{Portfolio web personal construido con Astro y Tailwind CSS, desplegado en Cloudflare. Implementa Edge Functions de Cloudflare para la redirección automática basada en el idioma del navegador, ofreciendo una experiencia multilingüe optimizada y de alto rendimiento.}
\resumeSubItem{MatricuApp}
{Solución a prueba técnica — una webapp de matrícula universitaria donde los estudiantes pueden iniciar sesión, matricular y retirar materias de su semestre correspondiente, con límite de créditos por semestre. Frontend construido con React y backend con Spring Boot, Maven, JPA, autenticación JWT con Spring Security y base de datos PostgreSQL en Neon.}
\resumeSubHeadingListEnd

%-----------Cursos y Certificaciones-----------------
\section{Cursos y Certificaciones}
\resumeSubHeadingListStart
\resumeSubItem{Universidad autónoma del caribe - Nuevas Tendencias de ingeniería}{Feb 2023}
\resumeSubItem{Udemy - Jetpack Compose desde cero, migra tus vistas de Android XML}{Abr 2024}
\resumeSubItem{Udemy - Curso Testing para Android con JUnit, Mockito, Espresso y TDD}{Feb 2024}
\resumeSubItem{Platzi - Curso Profesional de Git y GitHub}{Sep 2019}
\resumeSubItem{Platzi - Herramientas de AI para Developers}{Abr 2025}
\resumeSubItem{Platzi - Diseño de interfaces Android con Android Studio}{Jul 2020}
\resumeSubItem{Platzi - Curso de Flutter}{Mar 2021}
\resumeSubItem{Platzi - Curso Avanzado de Flutter}{Abr 2022}
\resumeSubItem{Platzi - Introducción a Java SE}{Oct 2020}
\resumeSubItem{Platzi - Java SE Orientado a Objetos}{Feb 2021}
\resumeSubItem{Platzi - Curso básico de Kotlin}{Oct 2019}
\resumeSubItem{Platzi - Curso de Kotlin para Android}{May 2020}
\resumeSubItem{Platzi - Curso de Introducción a Terminal y Linea de Comandos}{Abr 2020}
\resumeSubItem{Platzi - Curso de Bases técnicas de Android}{Nov 2019}
\resumeSubHeadingListEnd

%--------SKILLS------------
\section{Skills}
\resumeSubHeadingListStart
\item{
            \textbf{Lenguajes}{: Java, Kotlin, SQL, TypeScript, JavaScript, HTML, CSS, Dart, Python, LaTeX}

      }
\item {
            \textbf{Tecnologías}{: AWS, Android, Astro, React, Tailwind CSS, Cloudflare Edge Functions, Google Web Toolkit (GWT), Jetpack Compose, Flutter, Spring Boot, Kubernetes, Docker, Docker Compose, PostgreSQL, MySQL, Firebase, Git, GitHub Actions, Espresso, Mockito, JUnit, Test-Driven Development, Clean Architecture, SOLID, Design Patterns, RESTful APIs, Microservices, CI/CD, Agile, Scrum, Maven, Gradle, Make, Jira}
      }
\resumeSubHeadingListEnd


%-------------------------------------------
\end{document}
